% Ad Atticum 10.7 --- headnote
% ad-atticum-10-07, c. 23 April 49 BC, Cumae

\textit{Cicero to Atticus, written from the Cuman villa
about the ninth day before the Kalends of May 49 BC
(the manuscript dateline: \textit{Scr.\ in Cumano circ.\
ix K.\ Mai.\ a.\ 705 (49)}). Caesar is in Spain; Pompey
has crossed to Greece; Atticus has signalled an
intention to retire to Apulia and Sipontum and to keep
out of the fight \textit{tergiversando} --- by quiet
evasion --- and Cicero opens with explicit endorsement
of that plan, while marking why his own case is not the
same. The political diagnosis is the sharpest he has
yet committed to writing: this is no longer a contest
about the commonwealth but a contest about kingship,
\textit{regnandi contentio}, in which the more modest
and upright \textit{rex} has been driven out, but in
which even the better claimant, if he wins, will win in
Sulla's manner and after Sulla's example.}

\textit{The middle section turns to Servius Sulpicius
Rufus: a hope has been raised that Sulpicius wants a
conference with him, and Cicero has dispatched his
freedman Philotimus with a letter. ``If he means to be
a man, an admirable \textit{[Greek: synodia]}'' --- a
travelling-companionship --- ``but if not, we shall be
the men we have always been.'' Section 3 reports on a
recent visit from Curio (whose judgement is that Caesar
is faltering on account of his unpopularity and is
worried about Sicily once Pompey begins to move) and
returns again to young Quintus: greed, the expectation
of a Caesarian bounty, but not, Cicero hopes, the
worse crime he had feared. The closing line --- ``we
shall regard our Epirus as ours still; but it had
seemed we were going to take other courses'' ---
glances at the estate at Buthrotum that Atticus still
holds open as a possible refuge, while marking that
the planned destination is now somewhere else.}
